% Migrated from D:\Latex\RabbitMQ_JS.tex -- preamble now lives in shared/notebooks.sty.
\documentclass[11pt,a4paper]{article}
\usepackage{notebooks}

\notebooktitle{RabbitMQ}
\notebooksubtitle{JavaScript}
\notebooktagline{Messaging patterns with amqplib}
\notebookauthor{Bhuvnesh Verma}
\notebooksource{https://www.rabbitmq.com/tutorials}

\begin{document}
\notebookcover

\section{Getting Started}

\field{Short description:}{RabbitMQ is a message broker. It accepts, stores, and forwards
messages between applications. Think of it as a post office: you drop mail into a post
box, and the postal service delivers it.}

\field{Theory:}{Applications do not talk to each other directly. A producer hands a message
to the broker, the broker stores it in a queue, and a consumer picks it up when it is
ready. This decouples the two sides in both time and space.}

\field{How it works:}{A producer publishes to an exchange. The exchange applies its routing
rules and places the message into one or more queues. A consumer subscribes to a queue and
receives messages as they arrive.}

\field{Use case example:}{An online store accepts an order and must send a confirmation
email, update inventory, and notify shipping. Publishing one order event lets three
independent services react without the store waiting for any of them.}

\subsection*{Key Points}

\begin{itemize}\itemsep2pt
  \item \textbf{Producer} sends messages
  \item \textbf{Queue} stores messages, like a post box
  \item \textbf{Consumer} receives messages
  \item \textbf{Exchange} routes messages to queues
  \item The Node.js client library is \lstinline|amqplib|
\end{itemize}

\subsection*{Notes}

\begin{itemize}\itemsep2pt
  \item Every example assumes a broker running on \lstinline|localhost| at port 5672
  \item \lstinline|amqplib| is promise based, so every call happens inside an async function
  \item A channel, not a connection, is the unit you publish and consume on
\end{itemize}

\newpage
\subsection{Installation and connection}


\begin{sh}
npm install amqplib
\end{sh}

\begin{js}
const amqp = require('amqplib');

async function main() {
    // 1. Open a TCP connection to the broker
    const connection = await amqp.connect('amqp://localhost');

    // 2. Open a channel: all AMQP operations happen here
    const channel = await connection.createChannel();

    // ... declare queues, publish, consume ...

    // 3. Always close when finished
    await channel.close();
    await connection.close();
}

main().catch(console.error);
\end{js}

\newpage
% =========================================================================
\section{Hello World}

\field{Short description:}{The simplest pattern. One producer sends one message to one
queue, and one consumer receives it.}

\field{Theory:}{Point to point messaging. Producer $\rightarrow$ Queue $\rightarrow$
Consumer. No exchange is used, so messages go through the default exchange.}

\field{How it works:}{The producer declares a queue and sends a message to it. The consumer
declares the same queue and consumes from it. Declaring is idempotent, so whichever side
starts first creates the queue.}

\field{Use case example:}{A simple notification service where an order placement service
sends a ``new order'' message to a queue, and an inventory service picks it up to update
stock.}

\subsection{Architecture diagram}
\mmd{01-hello-world.png}

\subsection*{Key Points}

\begin{itemize}\itemsep2pt
  \item The queue must be declared before sending or receiving
  \item An empty exchange name (\lstinline|''|) uses the \textbf{default exchange}
  \item \lstinline|noAck: true| means the broker treats a message as delivered immediately
\end{itemize}

\subsection*{Notes}

\begin{itemize}\itemsep2pt
  \item Run the producer once; the consumer stays running
  \item Messages are \textbf{not persistent} by default and are lost if RabbitMQ restarts
\end{itemize}

\newpage
\subsection{Code examples}

\subsubsection*{Producer (\texttt{send.js})}
\begin{js}
const amqp = require('amqplib');

async function send() {
    const connection = await amqp.connect('amqp://localhost');
    const channel = await connection.createChannel();

    const queue = 'hello';
    await channel.assertQueue(queue, { durable: false });

    const msg = 'Hello World!';
    channel.sendToQueue(queue, Buffer.from(msg));
    console.log(" [x] Sent '%s'", msg);

    setTimeout(() => { connection.close(); }, 500);
}
send();
\end{js}

\subsubsection*{Consumer (\texttt{receive.js})}
\begin{js}
const amqp = require('amqplib');

async function receive() {
    const connection = await amqp.connect('amqp://localhost');
    const channel = await connection.createChannel();

    const queue = 'hello';
    await channel.assertQueue(queue, { durable: false });

    console.log(" [*] Waiting for messages. To exit press CTRL+C");

    channel.consume(queue, (msg) => {
        console.log(" [x] Received '%s'", msg.content.toString());
    }, { noAck: true });
}
receive();
\end{js}

\newpage
% =========================================================================
\section{Work Queues}

\field{Short description:}{Distributes time consuming tasks among multiple workers. This is
the competing consumers pattern.}

\field{Theory:}{Avoid doing resource intensive work inside a request. Schedule it to be done
later by a pool of workers that share one queue.}

\field{How it works:}{The producer sends tasks to a queue. Multiple workers consume from the
same queue, and tasks are handed out \textbf{round robin}. Each worker acknowledges a task
only once it has finished, so nothing is lost if a worker dies mid job.}

\field{Use case example:}{A video processing service receives upload requests. Each video is
transcoded by a worker. Workers run as multiple instances and tasks are distributed evenly.
If a worker crashes, its task is re-queued for another worker.}

\subsection{Architecture diagram}
\mmd{02-work-queues.png}

\subsection*{Key Points}

\begin{itemize}\itemsep2pt
  \item \lstinline|durable: true| makes the queue survive a broker restart
  \item \lstinline|persistent: true| makes the message survive a broker restart
  \item \lstinline|prefetch(1)| gives fair dispatch, so no worker is handed a backlog
  \item \textbf{Always acknowledge}, otherwise messages can be lost
\end{itemize}

\subsection*{Notes}

\begin{itemize}\itemsep2pt
  \item Round robin distribution gives each worker one message at a time
  \item If a worker dies without acknowledging, the message is re-queued
  \item Durability and persistence are separate settings, and you usually want both
\end{itemize}

\newpage
\subsection{Code examples}

\subsubsection*{Producer (\texttt{new\_task.js})}
\begin{js}
const amqp = require('amqplib');

async function send() {
    const connection = await amqp.connect('amqp://localhost');
    const channel = await connection.createChannel();

    const queue = 'task_queue';
    await channel.assertQueue(queue, { durable: true });

    const msg = process.argv.slice(2).join(' ') || "Hello World!";
    channel.sendToQueue(queue, Buffer.from(msg), { persistent: true });
    console.log(" [x] Sent '%s'", msg);

    setTimeout(() => { connection.close(); }, 500);
}
send();
\end{js}

\subsubsection*{Worker (\texttt{worker.js})}
\begin{js}
const amqp = require('amqplib');

async function receive() {
    const connection = await amqp.connect('amqp://localhost');
    const channel = await connection.createChannel();

    const queue = 'task_queue';
    await channel.assertQueue(queue, { durable: true });

    channel.prefetch(1);
    console.log(" [*] Waiting for messages. To exit press CTRL+C");

    channel.consume(queue, (msg) => {
        const secs = msg.content.toString().split('.').length - 1;
        console.log(" [x] Received %s", msg.content.toString());
        setTimeout(() => {
            console.log(" [x] Done");
            channel.ack(msg);
        }, secs * 1000);
    }, { noAck: false });
}
receive();
\end{js}

\newpage
% =========================================================================
\section{Publish and Subscribe}

\field{Short description:}{Broadcasts every message to \textbf{multiple consumers} at once.}

\field{Theory:}{The producer sends to a \textbf{fanout exchange}. The exchange forwards each
message to \textbf{all bound queues}, and every consumer owns its own queue.}

\field{How it works:}{The producer declares a fanout exchange. Each consumer creates a
temporary, exclusive queue and binds it to that exchange. Every message published to the
exchange is copied into all bound queues.}

\field{Use case example:}{A logging system where every log line must reach several sinks at
once: a file logger, a monitoring dashboard, and an alerting service. Each sink runs as a
separate consumer with its own queue.}

\subsection{Architecture diagram}
\mmd{03-pub-sub.png}

\subsection*{Key Points}

\begin{itemize}\itemsep2pt
  \item A \textbf{fanout exchange} ignores routing keys and broadcasts to everything bound
  \item Temporary queues declared as \lstinline|''| get a random generated name
  \item \lstinline|exclusive: true| deletes the queue when its consumer disconnects
  \item Each consumer receives \textbf{all} messages
\end{itemize}

\subsection*{Notes}

\begin{itemize}\itemsep2pt
  \item This differs from Work Queues, where consumers share the messages
  \item Here each consumer gets its own \textbf{copy} of every message
  \item Messages published while no queue is bound are simply discarded
\end{itemize}

\newpage
\subsection{Code examples}

\subsubsection*{Producer (\texttt{emit\_log.js})}
\begin{js}
const amqp = require('amqplib');

async function send() {
    const connection = await amqp.connect('amqp://localhost');
    const channel = await connection.createChannel();

    const exchange = 'logs';
    await channel.assertExchange(exchange, 'fanout', { durable: false });

    const msg = process.argv.slice(2).join(' ') || "info: Hello World!";
    channel.publish(exchange, '', Buffer.from(msg));
    console.log(" [x] Sent '%s'", msg);

    setTimeout(() => { connection.close(); }, 500);
}
send();
\end{js}

\subsubsection*{Consumer (\texttt{receive\_logs.js})}
\begin{js}
const amqp = require('amqplib');

async function receive() {
    const connection = await amqp.connect('amqp://localhost');
    const channel = await connection.createChannel();

    const exchange = 'logs';
    await channel.assertExchange(exchange, 'fanout', { durable: false });

    const q = await channel.assertQueue('', { exclusive: true });
    await channel.bindQueue(q.queue, exchange, '');

    console.log(" [*] Waiting for logs. To exit press CTRL+C");

    channel.consume(q.queue, (msg) => {
        console.log(" [x] %s", msg.content.toString());
    }, { noAck: true });
}
receive();
\end{js}

\newpage
% =========================================================================
\section{Routing}

\field{Short description:}{Selectively receives messages based on \textbf{routing keys},
matched exactly.}

\field{Theory:}{Uses a \textbf{direct exchange}. A message is routed to every queue whose
\textbf{binding key exactly matches} the message routing key.}

\field{How it works:}{The producer publishes with a routing key such as ``error'' or
``info''. Consumers bind their queues to the exchange with the specific keys they care
about, and the exchange delivers only to the queues that match.}

\field{Use case example:}{A log aggregator where services emit logs at different severity
levels. Only ``error'' logs reach the on call pager, while ``info'' and ``warning'' go to
database storage. Each consumer binds to the severities it wants.}

\subsection{Architecture diagram}
\mmd{04-routing.png}

\subsection*{Key Points}

\begin{itemize}\itemsep2pt
  \item A \textbf{direct exchange} routes by exact match
  \item Several queues may bind with different routing keys
  \item One queue may bind with \textbf{multiple} routing keys
\end{itemize}

\subsection*{Notes}

\begin{itemize}\itemsep2pt
  \item Run several consumers with different severity arguments to see the split
  \item Example: \lstinline|node receive_logs_direct.js error warning|
  \item A message whose key matches nothing is dropped unless you publish with
        \lstinline|mandatory|
\end{itemize}

\newpage
\subsection{Code examples}

\subsubsection*{Producer (\texttt{emit\_log\_direct.js})}
\begin{js}
const amqp = require('amqplib');

async function send() {
    const connection = await amqp.connect('amqp://localhost');
    const channel = await connection.createChannel();

    const exchange = 'direct_logs';
    await channel.assertExchange(exchange, 'direct', { durable: false });

    const severity = process.argv[2] || 'info';
    const msg = process.argv.slice(3).join(' ') || "Hello World!";
    channel.publish(exchange, severity, Buffer.from(msg));
    console.log(" [x] Sent %s: '%s'", severity, msg);

    setTimeout(() => { connection.close(); }, 500);
}
send();
\end{js}

\subsubsection*{Consumer (\texttt{receive\_logs\_direct.js})}
\begin{js}
const amqp = require('amqplib');

async function receive() {
    const connection = await amqp.connect('amqp://localhost');
    const channel = await connection.createChannel();

    const exchange = 'direct_logs';
    await channel.assertExchange(exchange, 'direct', { durable: false });

    const q = await channel.assertQueue('', { exclusive: true });

    const severities = process.argv.slice(2) || ['info'];
    for (const severity of severities) {
        await channel.bindQueue(q.queue, exchange, severity);
    }

    console.log(" [*] Waiting for %s. To exit press CTRL+C", severities);

    channel.consume(q.queue, (msg) => {
        console.log(" [x] %s: '%s'", msg.fields.routingKey, msg.content.toString());
    }, { noAck: true });
}
receive();
\end{js}

\newpage
% =========================================================================
\section{Topics}

\field{Short description:}{Receives messages by \textbf{pattern matching} on routing keys,
using wildcards.}

\field{Theory:}{Uses a \textbf{topic exchange}. Routing keys are dot separated words, and
consumers bind with patterns built from two wildcards: \lstinline|*| matches exactly one
word, \lstinline|#| matches zero or more words.}

\field{How it works:}{The producer publishes with a routing key such as
\lstinline|stock.usd.nyse|. A consumer binds a pattern such as \lstinline|*.usd.#| and
receives every message whose key fits that shape.}

\field{Use case example:}{A stock price feed. One client wants everything from the NYSE
(\lstinline|*.nyse.*|), another wants only USD denominated stocks from any exchange
(\lstinline|*.usd.*|), and a third wants the whole energy sector (\lstinline|energy.#|).}

\subsection{Architecture diagram}
\mmd{05-topics.png}

\subsection*{Key Points}

\begin{itemize}\itemsep2pt
  \item \lstinline|*| matches one word: \lstinline|*.stock.*| matches \lstinline|usd.stock.nyse|
  \item \lstinline|#| matches zero or more: \lstinline|#.error| matches \lstinline|app.service.error|
  \item Binding \lstinline|*.orange.*| matches \lstinline|quick.orange.rabbit|
  \item Binding \lstinline|lazy.#| matches both \lstinline|lazy.orange.rabbit| and \lstinline|lazy|
\end{itemize}

\subsection*{Notes}

\begin{itemize}\itemsep2pt
  \item A topic exchange bound with \lstinline|#| behaves exactly like a fanout exchange
  \item Bound with a key containing no wildcard, it behaves like a direct exchange
  \item Very flexible, and the usual choice for complex event filtering
\end{itemize}

\newpage
\subsection{Code examples}

\subsubsection*{Producer (\texttt{emit\_log\_topic.js})}
\begin{js}
const amqp = require('amqplib');

async function send() {
    const connection = await amqp.connect('amqp://localhost');
    const channel = await connection.createChannel();

    const exchange = 'topic_logs';
    await channel.assertExchange(exchange, 'topic', { durable: false });

    const routingKey = process.argv[2] || 'anonymous.info';
    const msg = process.argv.slice(3).join(' ') || "Hello World!";
    channel.publish(exchange, routingKey, Buffer.from(msg));
    console.log(" [x] Sent %s: '%s'", routingKey, msg);

    setTimeout(() => { connection.close(); }, 500);
}
send();
\end{js}

\subsubsection*{Consumer (\texttt{receive\_logs\_topic.js})}
\begin{js}
const amqp = require('amqplib');

async function receive() {
    const connection = await amqp.connect('amqp://localhost');
    const channel = await connection.createChannel();

    const exchange = 'topic_logs';
    await channel.assertExchange(exchange, 'topic', { durable: false });

    const q = await channel.assertQueue('', { exclusive: true });

    const bindingKeys = process.argv.slice(2) || ['#'];
    for (const key of bindingKeys) {
        await channel.bindQueue(q.queue, exchange, key);
    }

    console.log(" [*] Waiting for %s. To exit press CTRL+C", bindingKeys);

    channel.consume(q.queue, (msg) => {
        console.log(" [x] %s: '%s'", msg.fields.routingKey, msg.content.toString());
    }, { noAck: true });
}
receive();
\end{js}

\newpage
% =========================================================================
\section{Remote Procedure Call}

\field{Short description:}{The client sends a request and waits for a response, much like
calling a function that happens to live on another machine.}

\field{Theory:}{The client sends a request carrying a \textbf{callback queue} name and a
\textbf{correlation ID}. The server processes it and publishes the result to that callback
queue. The client listens on the callback queue and matches replies by correlation ID.}

\field{How it works:}{The client creates a temporary exclusive queue for replies and sets
\lstinline|replyTo| and \lstinline|correlationId| on the request. The server consumes the
request, computes, and sends the result to \lstinline|replyTo| with the same correlation ID.
The client resolves the matching promise.}

\field{Use case example:}{A web frontend needs a Fibonacci number that is expensive to
compute, so the work is offloaded to an RPC server. The frontend sends the input and waits
for the result to come back.}

\subsection{Architecture diagram}
\mmd{06-rpc.png}

\subsection*{Key Points}

\begin{itemize}\itemsep2pt
  \item The \textbf{callback queue} is where the server sends the response
  \item The \textbf{correlation ID} matches a response to its request and prevents mixups
  \item The client suspends until the reply arrives
\end{itemize}

\subsection*{Notes}

\begin{itemize}\itemsep2pt
  \item RPC is synchronous from the caller's point of view, so it adds latency
  \item Good when you genuinely need a result back, otherwise prefer a plain queue
  \item Always set a timeout, since a dead server means a promise that never resolves
\end{itemize}

\newpage
\subsection{Code examples}

\subsubsection*{RPC Server (\texttt{rpc\_server.js})}
\begin{js}[basicstyle=\ttfamily\scriptsize]
const amqp = require('amqplib');

async function fib(n) {
    return n < 2 ? n : await fib(n - 1) + await fib(n - 2);
}

async function startServer() {
    const connection = await amqp.connect('amqp://localhost');
    const channel = await connection.createChannel();

    const queue = 'rpc_queue';
    await channel.assertQueue(queue, { durable: false });
    channel.prefetch(1);
    console.log(' [x] Awaiting RPC requests');

    channel.consume(queue, async (msg) => {
        const n = parseInt(msg.content.toString());
        console.log(` [.] fib(${n})`);
        const result = await fib(n);
        channel.sendToQueue(msg.properties.replyTo, Buffer.from(result.toString()), {
            correlationId: msg.properties.correlationId
        });
        channel.ack(msg);
    });
}
startServer();
\end{js}

\subsubsection*{RPC Client (\texttt{rpc\_client.js})}
\begin{js}[basicstyle=\ttfamily\scriptsize]
const amqp = require('amqplib');
const { v4: uuidv4 } = require('uuid');

class RpcClient {
    async connect() {
        this.connection = await amqp.connect('amqp://localhost');
        this.channel = await this.connection.createChannel();
        const q = await this.channel.assertQueue('', { exclusive: true });
        this.callbackQueue = q.queue;
    }

    async call(n) {
        const corrId = uuidv4();
        return new Promise((resolve) => {
            this.channel.consume(this.callbackQueue, (msg) => {
                if (msg.properties.correlationId === corrId) {
                    resolve(parseInt(msg.content.toString()));
                }
            }, { noAck: true });

            this.channel.sendToQueue('rpc_queue', Buffer.from(n.toString()), {
                correlationId: corrId,
                replyTo: this.callbackQueue
            });
        });
    }
}

async function run() {
    const client = new RpcClient();
    await client.connect();
    console.log(" [x] Requesting fib(30)");
    console.log(" [.] Got %d", await client.call(30));
}
run();
\end{js}

\newpage
% =========================================================================
\section{Publisher Confirms}

\field{Short description:}{Guarantees that the broker \textbf{actually received} the
messages you published.}

\field{Theory:}{The channel is put into confirm mode. Every published message is then either
acknowledged or negatively acknowledged by the broker. For persistent messages the
confirmation arrives only \textbf{after} the message has been written to disk.}

\field{How it works:}{Create the channel with \lstinline|createConfirmChannel()|. After each
publish the broker replies with \lstinline|Basic.Ack| or \lstinline|Basic.Nack|. You can
wait for these synchronously with \lstinline|waitForConfirms()| or handle them
asynchronously through channel events.}

\field{Use case example:}{An order processing system that cannot lose an order. Each order
is published with confirms so the service knows the broker has accepted and persisted it
before moving on. If no confirm arrives within a timeout, the service retries.}

\subsection{Architecture diagram}
\mmd[0.30]{07-publisher-confirms.png}

\subsection*{Key Points}

\begin{itemize}\itemsep2pt
  \item \lstinline|createConfirmChannel()| enables confirms
  \item \lstinline|Basic.Ack| means the broker took responsibility for the message
  \item \lstinline|Basic.Nack| means it did not, for example a routing error
  \item \lstinline|waitForConfirms()| blocks until every outstanding message is confirmed
\end{itemize}

\subsection*{Notes}

\begin{itemize}\itemsep2pt
  \item Confirms say nothing about consumers, only about the broker
  \item Pair them with the \lstinline|mandatory| flag to catch unroutable messages
  \item Synchronous confirms are slow, so use the async form for high throughput
\end{itemize}

\newpage
\subsection{Code examples}

\subsubsection*{Synchronous confirms (\texttt{confirm\_publisher.js})}
\begin{js}[basicstyle=\ttfamily\scriptsize]
const amqp = require('amqplib');

async function publish() {
    const connection = await amqp.connect('amqp://localhost');
    const channel = await connection.createConfirmChannel();  // Confirm channel

    const queue = 'confirm_queue';
    await channel.assertQueue(queue, { durable: true });

    // Publish with a per-message callback
    channel.sendToQueue(queue, Buffer.from('Critical message'), { persistent: true },
        (err, ok) => {
            if (err) {
                console.log(" [x] Message NOT confirmed");
            } else {
                console.log(" [x] Message confirmed");
            }
        });

    await channel.waitForConfirms();
    console.log(" [x] All messages confirmed");

    setTimeout(() => { connection.close(); }, 500);
}
publish();
\end{js}

\subsubsection*{Asynchronous confirms (high throughput)}
\begin{js}[basicstyle=\ttfamily\scriptsize]
const amqp = require('amqplib');

async function publish() {
    const connection = await amqp.connect('amqp://localhost');
    const channel = await connection.createConfirmChannel();

    const queue = 'confirm_queue';
    await channel.assertQueue(queue, { durable: true });

    // Listen for confirms instead of waiting on each publish
    channel.on('ack', (msg) => console.log(' [✓] Message acked'));
    channel.on('nack', (msg) => console.log(' [✗] Message nacked'));

    for (let i = 0; i < 10; i++) {
        channel.sendToQueue(queue, Buffer.from(`Message ${i}`), { persistent: true });
        console.log(` [x] Sent message ${i}`);
    }

    await channel.waitForConfirms();
    console.log(" [x] All messages confirmed");

    setTimeout(() => { connection.close(); }, 500);
}
publish();
\end{js}

\newpage
% =========================================================================
\section{Queue Types}

\field{Short description:}{RabbitMQ offers three queue types. The type is fixed at
declaration time and decides how the queue stores messages and survives a node failure.}

\field{Theory:}{}

\begin{itemize}\itemsep2pt
  \item A \textbf{classic} queue lives on one node
  \item A \textbf{quorum} queue replicates its contents across an odd number of nodes using
        the Raft consensus algorithm
  \item A \textbf{stream} keeps an append only log that many consumers can read
        independently, each at its own offset
\end{itemize}

\field{How it works:}{You pick the type with the \lstinline|x-queue-type| argument when
declaring the queue. Classic is the default. Quorum trades a little throughput for
durability. A stream does not remove a message when it is read, so the same message can be
replayed by a new consumer later.}

\field{Use case example:}{Payment confirmations go to a quorum queue because losing one is
unacceptable. Thumbnail generation jobs go to a classic queue because they are cheap to
retry. An event feed that analytics, billing, and audit all read goes to a stream.}

\begin{table}[h]
\centering
\begin{tabularx}{\textwidth}{@{}lXX@{}}
\toprule
\rowcolor{tablehead}
\textbf{\color{accent}Type} & \textbf{\color{accent}Stores} & \textbf{\color{accent}Pick it when} \\
\midrule
Classic & One node, optionally on disk & Throughput matters more than surviving a node loss \\
Quorum  & Replicated by Raft, always on disk & The message must not be lost \\
Stream  & Append only log, replayable & Many consumers read the same history \\
\bottomrule
\end{tabularx}
\end{table}

\subsection*{Key Points}

\begin{itemize}\itemsep2pt
  \item The type cannot be changed later, you must delete and redeclare the queue
  \item Quorum queues need a majority of nodes alive, so use 3 or 5, never an even number
  \item A stream keeps messages after delivery, bounded by
        \lstinline|x-max-length-bytes| or \lstinline|x-max-age|
\end{itemize}

\subsection*{Notes}

\begin{itemize}\itemsep2pt
  \item Quorum queues do not support priorities or the \lstinline|exclusive| flag
  \item Mirrored classic queues are deprecated, quorum queues replace them
  \item Redeclaring an existing queue with different arguments raises
        \lstinline|PRECONDITION_FAILED| and kills the channel
\end{itemize}

\newpage
\subsection{Code examples}

\subsubsection*{Declaring each type (\texttt{queue\_types.js})}
\begin{js}
const amqp = require('amqplib');

async function declare() {
    const connection = await amqp.connect('amqp://localhost');
    const channel = await connection.createChannel();

    // Classic: the default, single node.
    await channel.assertQueue('jobs.classic', { durable: true });

    // Quorum: replicated across the cluster, always durable.
    await channel.assertQueue('payments.quorum', {
        durable: true,
        arguments: {
            'x-queue-type': 'quorum',
            'x-delivery-limit': 5   // give up after 5 redeliveries
        }
    });

    // Stream: append only log, capped by size and age.
    await channel.assertQueue('events.stream', {
        durable: true,
        arguments: {
            'x-queue-type': 'stream',
            'x-max-length-bytes': 2_000_000_000,
            'x-max-age': '7D'
        }
    });

    console.log(' [x] Declared classic, quorum and stream queues');
    await connection.close();
}
declare();
\end{js}

\subsubsection*{Reading a stream from an offset (\texttt{stream\_consumer.js})}
\begin{js}
const amqp = require('amqplib');

async function consume() {
    const connection = await amqp.connect('amqp://localhost');
    const channel = await connection.createChannel();

    // Streams require a prefetch, the broker refuses the consumer without one.
    await channel.prefetch(100);

    await channel.consume('events.stream', (msg) => {
        console.log(' [x] offset %s: %s',
            msg.properties.headers['x-stream-offset'],
            msg.content.toString());
        channel.ack(msg);
    }, {
        noAck: false,
        arguments: { 'x-stream-offset': 'first' }  // or 'last', 'next', a number
    });
}
consume();
\end{js}

\newpage
% =========================================================================
\section{Message Structure and Message Properties}

\field{Short description:}{Every message carries a body of raw bytes plus a set of
properties and headers that travel with it.}

\field{Theory:}{An AMQP message has three parts.}

\begin{itemize}\itemsep2pt
  \item The \textbf{body} is an opaque byte array, RabbitMQ never looks inside it
  \item The \textbf{properties} are a fixed set of fields defined by the protocol, such as
        \lstinline|contentType| and \lstinline|correlationId|
  \item The \textbf{headers} are a free form map you fill with anything your application
        needs
\end{itemize}

\field{How it works:}{In amqplib the body is always a \lstinline|Buffer|, so objects must
be serialised before publishing and parsed after receiving. Properties are passed as the
third argument to \lstinline|publish| and read back on
\lstinline|msg.properties|. Delivery metadata such as the routing key and the
delivery tag arrives separately on \lstinline|msg.fields|.}

\field{Use case example:}{A service publishes JSON with
\lstinline|contentType: 'application/json'|, a \lstinline|messageId| for deduplication, and
a \lstinline|traceId| header so the request can be followed across every service that
touches it.}

\begin{table}[h]
\centering
\begin{tabularx}{\textwidth}{@{}lX@{}}
\toprule
\rowcolor{tablehead}
\textbf{\color{accent}Property} & \textbf{\color{accent}What it is for} \\
\midrule
\lstinline|persistent|     & Write the message to disk so a broker restart keeps it \\
\lstinline|contentType|    & How the consumer should parse the body \\
\lstinline|correlationId|  & Ties an RPC reply back to its request \\
\lstinline|replyTo|        & Queue the reply should be sent to \\
\lstinline|messageId|      & Unique id, used for deduplication and logging \\
\lstinline|expiration|     & Per message TTL in milliseconds, as a string \\
\lstinline|priority|       & Rank inside a priority queue \\
\lstinline|headers|        & Your own key value map, also used by header exchanges \\
\bottomrule
\end{tabularx}
\end{table}

\subsection*{Key Points}

\begin{itemize}\itemsep2pt
  \item The body is bytes, the broker never parses or validates it
  \item \lstinline|persistent: true| plus a durable queue is what survives a restart
  \item \lstinline|expiration| is a string of milliseconds, not a number
\end{itemize}

\subsection*{Notes}

\begin{itemize}\itemsep2pt
  \item Keep the body small, put routing information in headers instead
  \item \lstinline|priority| needs \lstinline|x-max-priority| set on the queue
  \item \lstinline|msg.fields.redelivered| tells you this message was delivered before
\end{itemize}

\newpage
\subsection{Code examples}

\subsubsection*{Publishing with properties (\texttt{publish\_properties.js})}
\begin{js}[basicstyle=\ttfamily\scriptsize]
const amqp = require('amqplib');
const crypto = require('crypto');

async function send() {
    const connection = await amqp.connect('amqp://localhost');
    const channel = await connection.createChannel();

    await channel.assertQueue('orders', { durable: true });

    const payload = { orderId: 42, total: 199.5, currency: 'EUR' };

    channel.sendToQueue('orders', Buffer.from(JSON.stringify(payload)), {
        persistent: true,                    // survive a broker restart
        contentType: 'application/json',
        contentEncoding: 'utf-8',
        messageId: crypto.randomUUID(),
        correlationId: 'checkout-42',
        timestamp: Date.now(),
        appId: 'checkout-service',
        type: 'order.created',
        expiration: '60000',                 // milliseconds, as a string
        headers: {
            traceId: 'abc-123',
            attempt: 1,
            source: 'web'
        }
    });

    console.log(' [x] Sent order with properties');
    setTimeout(() => connection.close(), 500);
}
send();
\end{js}

\subsubsection*{Reading them back (\texttt{read\_properties.js})}
\begin{js}[basicstyle=\ttfamily\scriptsize]
const amqp = require('amqplib');

async function receive() {
    const connection = await amqp.connect('amqp://localhost');
    const channel = await connection.createChannel();

    await channel.assertQueue('orders', { durable: true });
    await channel.prefetch(1);

    channel.consume('orders', (msg) => {
        const { properties: p, fields: f, content } = msg;

        console.log(' [x] routingKey  :', f.routingKey);
        console.log('     redelivered :', f.redelivered);
        console.log('     messageId   :', p.messageId);
        console.log('     type        :', p.type);
        console.log('     traceId     :', p.headers.traceId);

        const body = p.contentType === 'application/json'
            ? JSON.parse(content.toString())
            : content.toString();

        console.log('     body        :', body);
        channel.ack(msg);
    }, { noAck: false });
}
receive();
\end{js}

\newpage
% =========================================================================
\section{Dead Letter Exchange (DLX) and TTL}

\field{Short description:}{Messages that expire, are rejected, or overflow a full queue are
republished to another exchange instead of being thrown away.}

\field{Theory:}{A queue declared with \lstinline|x-dead-letter-exchange| forwards any
\textbf{dead lettered} message to that exchange. A message is dead lettered when it is
nacked with \lstinline|requeue: false|, when its TTL expires, when the queue is full, or
when the delivery limit is reached.}

\field{How it works:}{TTL is set either on the queue with
\lstinline|x-message-ttl|, which applies to everything in it, or per message with the
\lstinline|expiration| property. When the timer runs out the broker removes the message and,
if a DLX is configured, republishes it there with an
\lstinline|x-death| header recording why and how many times.}

\field{Use case example:}{A retry pipeline. The main queue dead letters failures into a
wait queue with a 30 second TTL, whose own DLX points back at the main queue. The message
comes back automatically half a minute later, and after a few rounds it lands in a parking
queue a human can inspect.}

\subsection{Architecture diagram}
\mmd{08-dlx-ttl.png}

\subsection*{Key Points}

\begin{itemize}\itemsep2pt
  \item Four causes of dead lettering: rejected, expired, maxlen, delivery limit
  \item \lstinline|x-dead-letter-routing-key| overrides the original routing key
  \item Per message \lstinline|expiration| and queue \lstinline|x-message-ttl| both apply,
        the shorter one wins
\end{itemize}

\subsection*{Notes}

\begin{itemize}\itemsep2pt
  \item Only the message at the head of a queue expires, so a slow head delays the rest
  \item The \lstinline|x-death| header carries the retry count, use it to stop a loop
  \item Never point a queue's DLX back at itself without a counter, that is an infinite loop
\end{itemize}

\newpage
\subsection{Code examples}

\subsubsection*{Dead letter setup (\texttt{dlx\_setup.js})}
\begin{js}[basicstyle=\ttfamily\scriptsize]
const amqp = require('amqplib');

async function setup() {
    const connection = await amqp.connect('amqp://localhost');
    const channel = await connection.createChannel();

    // The exchange and queue that collect the failures.
    await channel.assertExchange('dlx', 'direct', { durable: true });
    await channel.assertQueue('orders.dead', { durable: true });
    await channel.bindQueue('orders.dead', 'dlx', 'orders');

    // The working queue: 30s TTL, capped length, the rest falls through to the DLX.
    await channel.assertQueue('orders', {
        durable: true,
        arguments: {
            'x-message-ttl': 30000,
            'x-max-length': 10000,
            'x-dead-letter-exchange': 'dlx',
            'x-dead-letter-routing-key': 'orders'
        }
    });

    await channel.prefetch(1);
    channel.consume('orders', (msg) => {
        try {
            handle(JSON.parse(msg.content.toString()));
            channel.ack(msg);
        } catch (err) {
            // requeue: false dead letters it, requeue: true would loop forever.
            channel.nack(msg, false, false);
        }
    }, { noAck: false });
}
setup();
\end{js}

\subsubsection*{Delayed retry with a wait queue (\texttt{retry.js})}
\begin{js}[basicstyle=\ttfamily\scriptsize]
const amqp = require('amqplib');
const MAX_RETRIES = 3;

async function setup() {
    const connection = await amqp.connect('amqp://localhost');
    const channel = await connection.createChannel();

    // Wait queue holds the message for 30s, then dead letters it back to the main queue.
    await channel.assertQueue('orders.wait', {
        durable: true,
        arguments: {
            'x-message-ttl': 30000,
            'x-dead-letter-exchange': '',
            'x-dead-letter-routing-key': 'orders'
        }
    });

    channel.consume('orders.dead', (msg) => {
        const deaths = msg.properties.headers['x-death'] || [];
        const count = deaths.length ? deaths[0].count : 0;
        if (count >= MAX_RETRIES) {
            channel.sendToQueue('orders.parked', msg.content, { persistent: true });
        } else {
            channel.sendToQueue('orders.wait', msg.content, { persistent: true });
        }
        channel.ack(msg);
    }, { noAck: false });
}
setup();
\end{js}

\newpage
% =========================================================================
\section{Mirrored and Replicated Queues}

\field{Short description:}{Keeping copies of a queue on several nodes so the messages
survive losing one of them.}

\field{Theory:}{Old \textbf{mirrored} classic queues copied a master queue to mirrors using
a policy, and are deprecated because a network partition could silently lose
acknowledged messages. \textbf{Quorum} queues replace them. They use Raft, where one
\textbf{leader} accepts every write and only confirms it once a \textbf{majority} of
replicas has it on disk.}

\field{How it works:}{Clients always talk to the leader, no matter which node they connect
to, and the broker forwards the traffic. If the leader dies, the remaining replicas elect a
new one from those with the most complete log, and clients reconnect to it without losing
confirmed messages.}

\field{Use case example:}{A three node cluster running payment events. Node one holds the
leader and is taken down for a kernel upgrade. Node two is elected leader, publishers
reconnect, and no confirmed payment is lost.}

\subsection{Architecture diagram}
\mmd{09-replicated-queues.png}

\subsection*{Key Points}

\begin{itemize}\itemsep2pt
  \item A quorum queue of 3 nodes tolerates 1 failure, 5 nodes tolerate 2
  \item Use an odd node count, an even one gives no extra safety
  \item Replication only helps together with publisher confirms and consumer acks
\end{itemize}

\subsection*{Notes}

\begin{itemize}\itemsep2pt
  \item Classic queue mirroring policies (\lstinline|ha-mode|) are removed in RabbitMQ 4
  \item Set \lstinline|x-quorum-initial-group-size| to control how many replicas start with
  \item Replicating every queue is wasteful, replicate only what must not be lost
\end{itemize}

\newpage
\subsection{Code examples}

\subsubsection*{Declaring a replicated queue (\texttt{quorum\_queue.js})}
\begin{js}[basicstyle=\ttfamily\scriptsize]
const amqp = require('amqplib');

async function setup() {
    // List every node so the client can fail over when one goes down.
    const connection = await amqp.connect([
        'amqp://localhost:5672',
        'amqp://node2:5672',
        'amqp://node3:5672'
    ]);
    const channel = await connection.createConfirmChannel();

    await channel.assertQueue('payments', {
        durable: true,
        arguments: {
            'x-queue-type': 'quorum',
            'x-quorum-initial-group-size': 3,  // leader plus two followers
            'x-delivery-limit': 5
        }
    });

    // The confirm only arrives once a majority of replicas has the message on disk.
    channel.sendToQueue('payments', Buffer.from('payment-42'), { persistent: true },
        (err) => {
            console.log(err ? ' [!] Rejected' : ' [x] Replicated and confirmed');
            connection.close();
        });
}

setup().catch(console.error);
\end{js}

\subsubsection*{Cluster and queue inspection}
\begin{sh}[basicstyle=\ttfamily\scriptsize]
# Cluster health and which nodes are running
rabbitmq-diagnostics cluster_status
rabbitmq-diagnostics check_running

# Queue type, leader node and the current replica set
rabbitmqctl list_queues name type leader members messages

# Move a leader off a node before taking it down for maintenance
rabbitmq-queues quorum_status payments
rabbitmq-upgrade drain

# Grow or shrink the replica set of an existing quorum queue
rabbitmq-queues add_member payments rabbit@node4
rabbitmq-queues delete_member payments rabbit@node2
\end{sh}

\newpage
% =========================================================================
\section{Arguments}

\field{Short description:}{Extra options passed when declaring a queue or starting a
consumer. They change the behaviour of the broker, unlike properties, which travel with a
single message.}

\field{Theory:}{Arguments come in three places.}

\begin{itemize}\itemsep2pt
  \item \textbf{Queue arguments} go in the \lstinline|arguments| map of
        \lstinline|assertQueue| and are fixed for the life of the queue
  \item \textbf{Consumer arguments} go in the \lstinline|arguments| map of
        \lstinline|consume| and affect only that consumer
  \item \textbf{Options} such as \lstinline|durable| or \lstinline|persistent| are plain
        fields, not \lstinline|x-| arguments
\end{itemize}

\begin{table}[h]
\centering
\small
\begin{tabularx}{\textwidth}{@{}llX@{}}
\toprule
\rowcolor{tablehead}
\textbf{\color{accent}Argument} & \textbf{\color{accent}Example} & \textbf{\color{accent}What it does} \\
\midrule
\lstinline|x-queue-type|              & \lstinline|'quorum'|         & Picks classic, quorum or stream \\
\lstinline|x-message-ttl|             & \lstinline|30000|            & Every message expires after 30 seconds \\
\lstinline|x-expires|                 & \lstinline|1800000|          & Delete the queue itself after 30 idle minutes \\
\lstinline|x-max-length|              & \lstinline|10000|            & Keep at most 10000 messages \\
\lstinline|x-max-length-bytes|        & \lstinline|1048576|          & Keep at most 1 MB of message bodies \\
\lstinline|x-overflow|                & \lstinline|'reject-publish'| & What a full queue does: \lstinline|drop-head| (default), \lstinline|reject-publish|, \lstinline|reject-publish-dlx| \\
\lstinline|x-dead-letter-exchange|    & \lstinline|'dlx'|            & Where rejected or expired messages are republished \\
\lstinline|x-dead-letter-routing-key| & \lstinline|'orders'|         & Routing key used when dead lettering \\
\lstinline|x-max-priority|            & \lstinline|10|               & Turns it into a priority queue with 10 levels \\
\lstinline|x-single-active-consumer|  & \lstinline|true|             & Only one consumer receives at a time, the rest wait \\
\lstinline|x-delivery-limit|          & \lstinline|5|                & Quorum only: dead letter after 5 redeliveries \\
\bottomrule
\end{tabularx}
\end{table}

\subsection*{Notes}

\begin{itemize}\itemsep2pt
  \item Arguments are fixed at declaration, changing one means deleting the queue or
        applying a \textbf{policy} instead
  \item An unknown \lstinline|x-| argument is ignored silently, so a typo simply does nothing
  \item Values must have the right type, \lstinline|'30000'| as a string is not the same as
        \lstinline|30000|
\end{itemize}

\newpage
\subsection{Code examples}

\subsubsection*{Consumer arguments}

\begin{table}[h]
\centering
\small
\begin{tabularx}{\textwidth}{@{}llX@{}}
\toprule
\rowcolor{tablehead}
\textbf{\color{accent}Argument} & \textbf{\color{accent}Example} & \textbf{\color{accent}What it does} \\
\midrule
\lstinline|x-priority|       & \lstinline|10|      & This consumer is served before lower priority ones \\
\lstinline|x-stream-offset|  & \lstinline|'first'| & Where to start reading a stream: \lstinline|first|, \lstinline|last|, \lstinline|next|, an offset or a timestamp \\
\lstinline|x-cancel-on-ha-failover| & \lstinline|true| & Cancel the consumer when the queue moves node \\
\bottomrule
\end{tabularx}
\end{table}

\subsubsection*{Arguments in practice (\texttt{arguments.js})}
\begin{js}[basicstyle=\ttfamily\scriptsize]
const amqp = require('amqplib');

async function setup() {
    const connection = await amqp.connect('amqp://localhost');
    const channel = await connection.createChannel();

    // durable is an option, everything under arguments is a broker behaviour.
    await channel.assertQueue('orders', {
        durable: true,
        arguments: {
            'x-queue-type': 'quorum',
            'x-message-ttl': 30000,
            'x-max-length': 10000,
            'x-overflow': 'reject-publish',      // fail the publish instead of dropping
            'x-dead-letter-exchange': 'dlx',
            'x-delivery-limit': 5
        }
    });

    // A backup consumer: it only receives when no higher priority one is connected.
    await channel.consume('orders', (msg) => {
        console.log(' [x] backup handled', msg.content.toString());
        channel.ack(msg);
    }, { noAck: false, arguments: { 'x-priority': 1 } });
}

setup().catch(console.error);
\end{js}

\newpage
% =========================================================================
\section{Quick Reference}

\begin{table}[h]
\centering
\begin{tabularx}{\textwidth}{@{}llX@{}}
\toprule
\rowcolor{tablehead}
\textbf{\color{accent}Pattern} & \textbf{\color{accent}Exchange} & \textbf{\color{accent}Use it when} \\
\midrule
Hello World        & default (\lstinline|''|) & One producer, one consumer, no routing \\
Work Queues        & default (\lstinline|''|) & Spread heavy tasks over a pool of workers \\
Publish/Subscribe  & fanout                   & Every consumer needs every message \\
Routing            & direct                   & Consumers pick messages by exact key \\
Topics             & topic                    & Consumers pick messages by key pattern \\
RPC                & default (\lstinline|''|) & You need a reply back from the worker \\
Publisher Confirms & any                      & Losing a message is unacceptable \\
\bottomrule
\end{tabularx}
\end{table}

\subsection*{Common practices}

\begin{itemize}\itemsep2pt
  \item Reuse one connection per process and open a channel per task
  \item Always handle \lstinline|connection.on('error')| and \lstinline|'close'|, then reconnect
  \item Set \lstinline|prefetch| on every consumer, otherwise one worker hoards the queue
  \item Acknowledge only after the work is genuinely finished
  \item Make queues durable and messages persistent when the data matters
  \item Never block the event loop inside a consumer callback
\end{itemize}


\end{document}
